\input{../../v3_rp/preamble}
\graphicspath{{../../}}
\makeatletter\def\input@path{{./}{../../}{../../v3_rp/}}\makeatother
\title{\bfseries When Does It Pay to Be Early?\\
\large Evidence from Every US Trademark Filing}
\author{Eric Silver}
\date{Brief, September 2026}

\begin{document}
\maketitle

\section*{The question}

Strategy offers two confident and opposite answers to whether a firm
gains by bringing a new kind of offering to market before others do.
The first says it does. Firms that create uncontested market space,
rather than competing in existing categories, grow faster and earn
more \citep{KimMauborgne2004}; early entrants secure scarce resources,
customer loyalty and scale \citep{LiebermanMontgomery1988}; and when a
new industry's field of entrants is later winnowed, it is the early
entrants who survive the shakeout and come to lead it
\citep{KlepperMiller1995, Klepper1997}. On this view, even if later
entry, mergers and failures thin a new market, the firms that opened it
should be over-represented among its survivors and among the companies
that reach public markets.

The second says it does not. Pioneers pay to prove that a market
exists, and followers enter behind them at lower cost and risk.
Where an innovation is easy to imitate and the pioneer lacks the
complementary assets to hold its position, the value it creates is
captured by imitators \citep{Teece1986}; histories of named markets
show that the firms remembered as pioneers were usually not the first
entrants, because the true first entrants mostly failed \citep{GolderTellis1993}; and in management
consulting, where service-mark filings trace which firms copy which new
offerings, rivals' new offerings are imitated routinely, the more so
when the innovating firm is prominent \citep{SemadeniAnderson2010}.
A third position holds that neither answer is general: the advantage
of being first depends on how quickly the technology and the market
are changing \citep{SuarezLanzolla2007}.

These claims have been tested mostly on named markets and on the firms
that survived long enough to be studied. The question here is what
happens across the whole population of entrants, including the ones
that failed: \emph{does being early with a new kind of offering pay,
and when does it not?}

\section*{The measure}

A companion paper \citep{SilverCorpus2026} re-parses every US trademark
filing into dated events and scores each filing's goods/services
description on \emph{lead}: whether the firms that filed in the same
industry over the following five years described their offerings more
like this filing than the firms of the five years before had, so that
a leading filing named a kind of offering that became more common after
it was filed, and a lagging one named a kind that was becoming less
common.

That reading of lead can be checked directly. Every description is
scored as a mix of fifty themes, and the themes' shares of an
industry's filings shift over time, some growing and some shrinking.
Among the most leading fifth of 2002--2018 registrations, 86\% named a
main theme that took a larger share of their industry's filings in the
five years after they filed than in the five years before; among the
most lagging fifth, 15\% did. Leading filings entered offerings that
others then crowded into.

\section*{What each side predicts}

\begin{enumerate}
\item[\textbf{H1}] \emph{Early pays.} Leading filings' products survive
longer, and their owners are more likely to raise private money, reach
public reporting and list. The early entrants who survive a shakeout
should be over-represented among the largest firms.
\item[\textbf{H2}] \emph{Early costs.} Leading filings' products fail
more often, and their owners are no more likely, or less likely, to
reach the capital markets, because the crowding that makes a filing
leading is the competition that erodes its position.
\item[\textbf{H3}] \emph{It depends.} The sign of the lead effect varies
with the pace of change, by era and by industry.
\end{enumerate}

\noindent The record supplies the outcomes. Product survival is the
sworn declaration that a registered mark's goods are still being sold,
due in the fifth to sixth year; the capital markets are a priced
private round after the owner's first filing, appearance in SEC
reporting, and an initial public offering.

\section*{The answer: being early costs, at every stage after registration}

\emph{Products described early fail more often.} Among 3.10 million
registrations issued 2002--2018, the most leading fifth were cancelled
at the five-year proof 1.4 percentage points more often than the most
lagging fifth (53.3\% against 51.9\%), within industry and year and
with drafting, counsel and filer characteristics held ($t = 8.3$). The
penalty is positive in three of every four industries and under every
construction of the measure tried, and it persists, at about half the
size, at the year-ten renewal.

\emph{Early firms reach the capital markets less often, not more.}
Among owners whose first registered filing was made 2009--2018, the
most leading fifth raised a priced round after that filing no more
often than the most lagging fifth (1.69\% against 1.74\%), reached SEC
reporting less often (0.36\% against 0.55\%), and listed less often
(0.17\% against 0.22\%). The owners likeliest to report and to list
are those whose first description used language their industry was
leaving behind --- mostly established firms filing a first trademark
late.

\emph{New business models attract money but not listings.} Debuts
described in the vocabulary of artificial intelligence or blockchain
raised private rounds about five times as often as the average debut
(7.8\% and 8.0\% against 1.6\%), yet listed at the average rate (0.16\%
to 0.22\% against 0.17\%). The prediction that the new models' early
entrants would reach public markets in numbers does not hold on this
record: investors funded the new vocabulary, and the listings did not
follow.

These results favour H2 over H1. The effect is modest --- a leading
filing's product is about 3\% more likely to be gone within six years
--- so being early is a cost to be weighed, not a verdict.

\section*{When the rule breaks}

\emph{In the years after the dot-com crash, leading technology filings
survived more often.} For technology filings made 2000--2004, the most
leading fifth failed less often than the most lagging fifth. Two things
produced that reversal. The computer and information-services
vocabulary that had led the late 1990s had fallen out of the leading
fifth, and its marks failed at among the era's highest rates; excluding
that one theme restores a penalty. And within themes, leading filings
in goods and retail were genuinely rewarded: most strongly in online
retail, where the most leading fifth failed 8.3 points less often than
the most lagging, and also in electronic devices and in media goods.
Leading filings in research and scientific services were penalized in
the same years. The exception is consistent with H3 --- a period of
rapid change in how goods reached customers --- though this record
cannot separate it from the selection of which firms were still filing
after the crash.

\emph{Unusual is not the same as early.} Filings whose language is
unusual for their industry survive the proof more often, not less, and
filings that carry a theme established in another industry into a new
one survive more often than their industry's average. What the record
penalizes is not being different; it is being where the crowd is about
to go.

\emph{Joining a wave costs little; arriving first within it costs
nothing more.} Filings that joined a theme whose share of its industry
was surging failed only slightly more often than others, and within a
wave the order of arrival made no difference. The one corpus-wide
surge, online software in 1999, is the exception: its filings failed
8.6 points more often than their industry and year.

\emph{Among the few firms that grew very large, being early neither
helped nor hurt.} Of the 500 highest-revenue firms the record links to
SEC filings, 20\% debuted in the most leading fifth, the share chance
would give.

\section*{What the answer does not say}

Lead measures the position of a firm's language relative to its
industry's, not the novelty of its product, and crowding in language is
not crowding in market share. A filing can lead because others imitated
it, because all of them responded to the same outside change, or by
coincidence. Acquisitions are invisible: a mark retired after its owner
was acquired counts as a product that failed, and an acquisition is
not a listing. And the socially valuable part of early entry may be the
one the record cannot price: the language each industry eventually
adopted was first filed, disproportionately, by firms that did not
survive to use it \citep{KerrNanda2014}.

\input{bib_brief}
\end{document}
